\section{Présentation de l’environnement}
Ce stage de 2\textsuperscript{ème} année a été réalisé au sein de \textbf{RUHENNA}, une activité artisanale portée par \textbf{Sumeyra Solak} et centrée sur l’art du henné (prestations pour événements et produits associés). La structure étant de petite taille, l’enjeu principal était de livrer une solution \textbf{simple à utiliser} et \textbf{administrable} au quotidien.

\section{Sujet et objectifs du stage}
Le sujet du stage était de participer à la conception et au développement d’une plateforme web \textit{full-stack} permettant :
\begin{itemize}
    \item de valoriser la marque et les réalisations (galerie, blog) ;
    \item de rendre l’offre lisible (catalogue) ;
    \item de structurer les demandes (contact et réservation) ;
    \item de permettre une gestion autonome via un espace d’administration.
\end{itemize}

\section{Outils et moyens utilisés}
Le projet s’appuie sur une stack web courante : \textbf{React/Next.js} (App Router), \textbf{Prisma} et \textbf{PostgreSQL} (hébergée sur \textbf{Neon}), \textbf{Tailwind CSS}, \textbf{Cloudinary} pour les images, et un déploiement sur \textbf{Vercel}.

\section{Annonce du plan}
Le rapport présente d’abord le contexte et les choix (contraintes, outils, méthode, alternatives), puis détaille la solution mise en place, les contributions et les résultats observables (exemples de parcours, gestion admin, mise en production). Il se termine par une synthèse et des perspectives, avec des annexes techniques.
